\clearpage\begingroup\SourceLeaf{591}
\ArithmeticHead{36}{L}{PHILIPP MAENNCHEN, GAUSS ALS ZAHLENRECHNER.}
\renewcommand{\thefootnote}{\arabic{footnote})}\setcounter{footnote}{0}
\begin{center}VIII.\\[3pt]{\large\textls[80]{Zur Cyklotechnie.}}\end{center}
\noindent\hspace*{\parindent}Wir kommen nun zu den merkwürdigsten Tabellen, die \textsc{Gauss} aufgestellt
hat, und die sowohl in bezug auf ihren eigentlichen Zweck, als auch in bezug
auf die bei der Aufstellung angewandte Methode bisher ein Rätsel geblieben
sind. Es sind dies die \textit{Tafeln zur Cyklotechnie}\footnote{Werke II, S. 477.}.

Der ursprüngliche Zweck dieser Tafeln war wohl, wie schon der Name
andeutet, die Erleichterung, die sie für die genaue Berechnung der Bögen
gewähren, deren Cotangenten gegebene rationale Zahlen sind. Aber \textsc{Schering}
schreibt schon in seinen Bemerkungen zu diesen Tafeln\footnote{Ebenda, S. 499 und 500.}:

»Die hierauf hinzielenden Entwicklungen, die sich in dem handschriftlichen Nachlass finden, sind wenig ausgedehnt.«

Daraus lässt sich vermuten, dass sie \textsc{Gauss} zu dem erwähnten Zweck
wenig benutzt hat. Halten wir aber die Tatsache daneben, dass die ersten
Rechnungen für die Tafeln der Zeit der Ausarbeitung der \textit{Disquisitiones Arithmeticae} angehören, und dass in den Jahren 1846 und 1847 noch an der Sichtung, Ordnung und Vervollständigung derselben gearbeitet wurde, so ist der
Schluss wohl berechtigt, dass \textsc{Gauss} mit diesen Tafeln andere Pläne verfolgte.
Es wird jedoch zweckmässig sein, zunächst der Herstellung dieser Tafeln nachzugehen und erst dann Vermutungen über den Zweck zu äussern.

Bei der Herstellung der Tafeln hat \textsc{Gauss}, wie \textsc{Schering} bereits richtig
vermutet hat, besondere Kunstgriffe angewandt. \textsc{Schering} vermochte indessen
aus dem handschriftlichen Nachlass nur eine Regel aufzufinden, die aus drei
Zahlen der Tafel eine vierte zu finden lehrt, während es sich doch hauptsächlich darum handelt, aus zwei Zahlen eine dritte zu ermitteln. Wir wollen
nun versuchen, zu zeigen, in welcher Weise die Tafeln vermutlich entstanden
sind.

Die erste Tafel enthält die \textls[120]{zerlegbaren} $a^2+1$. Zerlegbar nennt \textsc{Gauss}
diejenigen Zahlen von der Form $a^2+1$, die nur Primteiler unter 200 haben.
Links stehen die Zahlen $a$, rechts die Primfaktoren von $a^2+1$, wobei noch
zu beachten ist, dass der Primfaktor 2 überall weggelassen wurde.
\clearpage\SourceLeaf{592}\ArithmeticHead{37}{R}{VIII. ZUR CYKLOTECHNIE.}
\setcounter{footnote}{0}
\begin{center}\renewcommand{\arraystretch}{1.13}\setlength{\tabcolsep}{5pt}
\begin{tabular}{r|l@{\qquad}r|l@{\qquad}r|l@{\qquad}r|l}
2 & $5$ & 19 & $181$ & 46 & $29\mathord{.}73$ & 93 & $5\mathord{.}5\mathord{.}173$\\
3 & $5$ & 21 & $17\mathord{.}13$ & 47 & $5\mathord{.}3\mathord{.}17$ & 98 & $5\mathord{.}17\mathord{.}113$\\
4 & $17$ & 22 & $5\mathord{.}97$ & 50 & $41\mathord{.}61$ & 99 & $13\mathord{.}13\mathord{.}29$\\
5 & $13$ & 23 & $5\mathord{.}53$ & 55 & $17\mathord{.}89$ & 100 & $37\mathord{.}137$\\
6 & $37$ & 27 & $5\mathord{.}73$ & 57 & $5\mathord{.}5\mathord{.}5\mathord{.}13$ & 105 & $37\mathord{.}149$\\
7 & $5\mathord{.}5$ & 28 & $5\mathord{.}157$ & 68 & $5\mathord{.}5\mathord{.}5\mathord{.}37$ & 111 & $61\mathord{.}101$\\
8 & $5\mathord{.}13$ & 30 & $17\mathord{.}53$ & 70 & $13^2\mathord{.}29$ & 112 & $5\mathord{.}13\mathord{.}193$\\
9 & $41$ & 31 & $13\mathord{.}37$ & 72 & $5\mathord{.}17\mathord{.}61$ & 117 & $5\mathord{.}37\mathord{.}37$\\
10 & $101$ & 32 & $5\mathord{.}5\mathord{.}41$ & 73 & $5\mathord{.}13\mathord{.}41$ & 119 & $73\mathord{.}97$\\
11 & $61$ & 33 & $5\mathord{.}109$ & 75 & $29\mathord{.}97$ & 123 & $5\mathord{.}17\mathord{.}89$\\
12 & $5\mathord{.}29$ & 34 & $13\mathord{.}89$ & 76 & $53\mathord{.}109$ & 128 & $5\mathord{.}29\mathord{.}113$\\
13 & $5\mathord{.}17$ & 37 & $5\mathord{.}137$ & 80 & $37\mathord{.}173$ & 129 & $53\mathord{.}157$\\
14 & $197$ & 38 & $5\mathord{.}17\mathord{.}17$ & 81 & $17\mathord{.}193$ & 132 & $5\mathord{.}5\mathord{.}17\mathord{.}41$\\
15 & $113$ & 41 & $29\mathord{.}29$ & 83 & $5\mathord{.}13\mathord{.}53$ & 133 & $5\mathord{.}29\mathord{.}61$\\
17 & $5\mathord{.}29$ & 43 & $5\mathord{.}5\mathord{.}37$ & 91 & $41\mathord{.}101$ & 142 & $5\mathord{.}37\mathord{.}109$\\
18 & $5\mathord{.}5\mathord{.}13$ & 44 & $13\mathord{.}149$ & & & &
\end{tabular}\end{center}
\noindent\hspace*{\parindent}Die Zahlen $a$, für die $a^2+1$ nicht zerlegbar ist, hat \textsc{Gauss} weggelassen.
— Im ersten Hunderter für $a$ empfiehlt sich ein Verfahren, das sicherlich
\textsc{Gauss} nicht entgangen ist, und das er bei einer andern Gelegenheit\footnote{Besprechung von \textsc{Burckhardt}, \textit{Tables des Diviseurs}, 1814, Werke II, S. 183.} warm
empfohlen hat, nämlich die Anwendung eines \textls[120]{Zahlensiebs}. Man sieht
nämlich leicht, dass zu 2, 7, 12, 17, 22 usw., also zu allen Zahlen von der
Form $5n+2$ der Faktor 5 gehört; ebenso gehört er zu allen Zahlen von der
Form $5n-2$. Der Faktor 13 findet sich zuerst bei 5; man darf ihn mit
Sicherheit erwarten bei $13n+5$ und bei $13n-5$. Der Faktor 17 findet sich
zuerst bei 4, folglich bei allen Zahlen von der Form $17n+4$ und $17n-4$\footnote{Nach diesem Verfahren hat \textsc{Euler} sämtliche $a^2+1$ bis $a=1500$ in Primfaktoren zerlegt: \textit{De numeris primis valde magnis.} Nov. Comm. \textbf{9} (1762/3), \textsc{Leonhardi Euleri}, Opera omnia, ser. I, vol. IV, S. 1—45. Es darf als sicher angenommen werden, dass \textsc{Gauss} diese Arbeit gekannt hat und durch sie beeinflusst worden ist.}.

Der Beweis ist einfach. Wenn $a^2+1$ den Primteiler $p$ hat, so hat, da
$(a\pm p)^2+1-(a^2+1)$ zweifellos immer durch $p$ teilbar ist, der Minuend
$(a\pm p)^2+1$ auch den Primteiler $p$, also gehört zu allen $np+a$ und $np-a$
der Teiler $p$. Weitere interessante Eigenschaften offenbaren sich auch bereits
auf dem obigen kleinen Anfangsstückchen der ersten Tabelle.
\clearpage\SourceLeaf{593}\ArithmeticHead{38}{L}{PHILIPP MAENNCHEN, GAUSS ALS ZAHLENRECHNER.}
\noindent\hspace*{\parindent}So wird man vor allen Dingen leicht entdecken, dass zu 7 dieselben Faktoren gehören, wie zu 2 und 3; zu 13 dieselben, wie zu 3 und 4; zu 21
dieselben, wie zu 4 und 5; \mbox{.\ .\ .} zu 111 dieselben, wie zu 10 und 11. Allgemein liefern $a^2+1$ und $(a+1)^2+1$ die Faktoren für $[a(a+1)+1]^2+1$.
Denn es ist
\[
\begin{aligned}
[a(a+1)+1]^2+1&=a^2(a+1)^2+2a(a+1)+2,\\
(a^2+1)\mathord{.}[(a+1)^2+1]&=a^2(a+1)^2+a^2+a^2+2a+1+1\\
&=a^2(a+1)^2+2a(a+1)+2,
\end{aligned}
\]
\noindent also
\[
[a(a+1)+1]^2+1=(a^2+1)\mathord{.}[(a+1)^2+1].
\]

Ebenso wird man leicht finden:

$a=3$ und $a=5$ liefern die Faktoren für $a=8$; 5 und 7 für 18; 7
und 9 für 32; 9 und 11 für 50, usw. Allgemein: $a$ (ungerade) und $a+2$
für $\frac{a(a+2)+1}{2}$. Denn
\[
(a^2+1)[(a+2)^2+1]=4\left[\left(\frac{a(a+2)+1}{2}\right)^2+1\right].
\]

Eine Beobachtung ähnlicher Art ist die folgende:

7 und 12 liefern die Faktoren für $\frac{7\mathord{.}12+1}{5}=17$, nur muss der Faktor 5
zweimal unterdrückt werden. Ebenso 8 und 13 für 21 mit derselben Forderung für den Faktor 5; 12 und 17 für $\frac{12\mathord{.}17+1}{5}=41$; allgemein: $a$ (nicht
durch 5 teilbar) und $a+5$ für $\frac{a(a+5)+1}{5}$. Denn es ist:
\[
(a^2+1)[(a+5)^2+1]=25\left[\left(\frac{a(a+5)+1}{5}\right)^2+1\right].
\]
\noindent Alle diese Sätze sind Sonderfälle der folgenden umfassenderen Formel:
\[
(a^2+1)[(a+p)^2+1]=[a(a+p)+1]^2+p^2.
\]
\noindent Wenn $a(a+p)+1$ durch $p$ teilbar ist (und dies ist der Fall, sobald zu $a$ der
Faktor $p$ gehört), dann ist die rechte Seite durch $p^2$ teilbar, und man erhält:
\[
(a^2+1)[(a+p)^2+1]=p^2\left[\left(\frac{a(a+p)+1}{p}\right)^2+1\right],
\]
\noindent eine Formel, in der sämtliche bisher besprochene Eigenschaften enthalten
sind. Ist dagegen $a(a+p)+1$ nicht durch $p$ teilbar, dann gehören die zu $a$
und $a+p$ gehörigen Faktoren in der Tabelle für $a^2+1$ zu der Zahl $a(a+p)+1$
\clearpage\SourceLeaf{594}\ArithmeticHead{39}{R}{VIII. ZUR CYKLOTECHNIE.}
\setcounter{footnote}{0}
\noindent \textls[120]{in der Tafel der} $a^2+p^2$. Z. B. zu 9 gehört 41, zu 12 gehört $5\mathord{.}29$, zu
$9\mathord{.}12+1=109$ gehört $5\mathord{.}29\mathord{.}41$ in der Tafel der zerlegbaren $a^2+3^2$.

Die bis jetzt gewonnenen Ergebnisse kann man zu mancherlei Rechenvorteilen benutzen, von denen hier nur einer angegeben werden soll: Haben
zwei Zahlen, die zu $a_1$ und $a_2$ der zerlegbaren $a^2+1$ gehören, den Faktor $\rho$
gemeinsam, so lässt $a_1a_2+1$, durch $\rho$ geteilt, entweder den Rest 0 oder den
Rest 2, und zwar den Rest 0, wenn $a_1-a_2$, den Rest 2, wenn $a_1+a_2$ durch
$\rho$ teilbar ist. Z. B. $233\mathord{.}1568+1$ muss durch 89 teilbar sein; denn $1568-233=1335=15\mathord{.}89$.
$743\mathord{.}8307+1$ lässt, durch 181 geteilt, den Rest 2;
denn $8307+743=9050=50\mathord{.}181$. Dass dies auch mancherlei praktische
Verwendbarkeit haben kann, liegt auf der Hand\footnote{Beweis: Haben $a_1^2+1$ und $a_2^2+1$ den Faktor $\rho$ gemeinsam, dann ist entweder $a_2=n\rho+a_1$ oder $a_2=n\rho-a_1$. Im ersten Fall ist $a_1a_2+1=a_1(n\rho+a_1)+1=a_1n\rho+(a_1^2+1)$; im zweiten Fall ist $a_1a_2-1=a_1(n\rho-a_1)-1=a_1n\rho-(a_1^2+1)$.}.

Wir wollen den Aufbau der Tafeln weiter fortführen. Wir haben
gesehen, dass man je nach der Wahl des Abstandes sowie des gemeinschaftlichen Faktors entweder in Tafel I (zerlegbare $a^2+1$) bleiben oder zu einer
anderen Tafel gelangen kann. Demnach kann man mit Tafel I bestimmte
Teile der übrigen Tafeln zusammenstellen. Aber auch umgekehrt kann man
von einer der folgenden Tafeln zur Tafel I zurückkommen. Ich führe zunächst ein Beispiel an:

Zu 63 in Tafel II gehört $29\mathord{.}137$; $63:2=31$, Rest 1; zu 31 gehört
$13\mathord{.}37$ in Tafel I; nun berechne ich $63\mathord{.}31+2=1955$, und hierzu gehört in
Tafel I $13\mathord{.}29\mathord{.}37\mathord{.}137$.

Ein anderes Beispiel:

Zu 96 in Tafel VII gehört $5\mathord{.}17\mathord{.}109$; $96:7=13$, Rest 5; zu 13 gehört $5\mathord{.}17$ in Tafel I; $13\mathord{.}96=1248$; $1248+7=1255$; $1255:5=251$; zu
251 gehört $17\mathord{.}17\mathord{.}109$ in Tafel I.

Damit kommen wir nun zu Verfahren, die \textsc{Gauss} \textls[120]{sicher angewandt
hat}, während dies bei den oben genannten Vorteilen nur höchst wahrscheinlich ist. Im \textsc{Gauss}schen Nachlass befindet sich ein Blatt, das ganz mit flüchtigen Zahlenrechnungen bedeckt ist, die auf den ersten Anblick wie einfache
Divisionen aussehen. So kam es wohl, dass dieses Blatt bisher nicht genügend
beachtet worden ist.
\clearpage\endgroup
